\section{Implementation and provenance details}
\label{app:impl}

\subsection{Structured network implementation}
\label{app:architecture}

The 229-dimensional canonical state has a five-value batch prefix followed by
16 padded 14-value event slots. Each slot contains an active flag, 12
normalized local features, and the stored prior action. The service index is
expanded to a 12-way one-hot code; selected/current status and relative slot
position are appended before encoding. One shared slot multilayer perceptron
maps each local vector to 128 features. The current slot, active-slot mean,
active-slot maximum, previously selected-slot mean, and future-slot mean form
five pooled views. A fusion multilayer perceptron combines these views with the
five-value prefix to produce 256 state features.

The actor maps the fused features through two 256-unit hidden layers to a gate
logit, latent normal mean, and log standard deviation. The
log standard deviation is clipped to $[-5,2]$. The initial gate bias gives a
positive-hold probability of approximately 0.35 and the conditional bias
favors short holds. The critic ensemble contains two independently
parameterized heads after one structured encoder; the target ensemble is an
exact structural copy. The value network has a separate encoder and scalar
head. Main modules use rectified linear unit activations and no layer
normalization.

\subsection{Frozen training values}
\label{app:training}

\Cref{tab:hparams} lists the principal final-profile values. All optimizers use
adaptive moment estimation (Adam). Stored rewards remain physical generalized passenger-seconds; numerical
conditioning is applied only inside the trainer. Actions, rewards,
continuations, durations, and terminal flags remain single precision.

\begin{table*}[tp]
\centering
\small
\caption{Final actor-training and architecture values.}
\label{tab:hparams}
\begin{tabular}{@{}p{0.54\linewidth}p{0.42\linewidth}@{}}
\toprule
Parameter & Value \\
\midrule
Canonical state width & 229 \\
Maximum event slots & 16 \\
Slot embedding and fused widths & 128 and 256, respectively \\
Actor/critic/value hidden widths & 256, 256 \\
Critic heads & 2 \\
Maximum hold & 60 s \\
Critic and simulator-actor minibatches & 256 and 128, respectively \\
Actor/critic/value/entropy learning rates & $3\times10^{-4}$ \\
Ratio learning rate & $10^{-4}$ \\
Target update rate & 0.005 \\
Simulator critic-minibatch fraction & 0.5 \\
IQL expectile & 0.7 \\
IQL exponential multiplier & 3 \\
Maximum IQL actor weight & 100 \\
Mixed-backup coefficient & 0.1 \\
Initial entropy temperature & 1.0 \\
Target entropy & 0.0 \\
Trainer reward scale $\kappa_r$ & $10^{-4}$ \\
Critic Huber threshold & 1.0 \\
Actor interpolation coefficient & $5\times10^{-5}$ \\
Offline IQL updates & 50,000 \\
Frozen-policy warmup & 5,000 events \\
Actor/temperature freeze endpoint & 10,000 events \\
Actor ramp & 10,000--15,000 events; multiplier approximately 0.1 to 1 \\
Ratio pretraining & 1,000 updates \\
Ratio update interval & 4 learner updates \\
Terminal actor taper & 15,000--25,000 events; multiplier linear 1 to 0.25 \\
Total learner updates & 70,000 \\
Fixed actor endpoint & 25,000 simulator events \\
Training seeds & 6201--6205 \\
\bottomrule
\end{tabular}
\end{table*}

The terminal-quarter branch reused each seed's sealed 10,000-event state and
continued without restarting optimizer, replay, or simulator chronology. All
15 terminal-taper branches passed artifact and training-stability checks; only
the preregistered quarter-rate branch supplied the five actors used by the
final candidate. The actor regression uses the weight in
\Cref{eq:iql_weight}; the online critic uses \Cref{eq:critic_target} with
\Cref{eq:mixed_backup,eq:sac_bootstrap}.

\subsection{Ratio diagnostic}
\label{app:ratio}

Target and simulator examples enter each domain-classifier loss in equal
numbers. Offline validation is split by complete environment block, and
simulator validation holds out complete trajectories. A diagnostic random
stream is separate from optimization sampling. \Cref{tab:ratio_diag} reports
the across-seed summary that determined interpretation of the ratio module.

\begin{table}[t]
\centering
\small
\caption{Domain-ratio diagnostics across five sealed actor-training seeds.}
\label{tab:ratio_diag}
\begin{tabular}{lrrr}
\toprule
Diagnostic & Mean & Minimum & Maximum \\
\midrule
Training joint AUC & 0.9477 & 0.9370 & 0.9580 \\
Held-out joint AUC & 0.4520 & 0.4366 & 0.4747 \\
Held-out raw ESS fraction & 0.0615 & 0.0100 & 0.1530 \\
\bottomrule
\end{tabular}
\end{table}

The separation between training and validation AUC, together with low raw
ESS, shows that the classifier did not supply a usable
generalizing correction. The optimizer used unit ratio weights, so the
diagnostic failure did not directly perturb critic updates.

\subsection{Dataset integrity and support}
\label{app:support}

The immutable replay contains 3,284,718 rows, 240 trajectories, 240 terminal
rows, and 1,025,206 zero-duration agent-environment-cycle rows. No nonfinite row
was found, and chronological passenger-ledger reconstruction had zero maximum
error. Exact zero hold accounts for 80.2246\% of actions, while 2,145 rows use
the exact 60-s upper endpoint. The audit groups uncertainty by complete
environment cells rather than treating rows as independent samples.

For the final low-load guard, 128,437 rows (3.91\%) satisfy the physical-state
conditions in the audited replay. Among these rows, 5.81\%, 3.69\%, and 2.17\%
have behavior holds of at least 20, 30, and 40 s, respectively. These values
motivate reporting the final 60-s action as an operational reserve outside
dense logged support.

\subsection{Seed-specific formal effects}
\label{app:seed_effect}

\Cref{tab:seed_effect} shows the candidate-minus-parent effect after averaging
each training seed over the ten formal blocks and three demand levels. The
direction is the same for all seeds.

\begin{table}[t]
\centering
\small
\caption{Formal candidate-minus-parent differences by training seed.}
\label{tab:seed_effect}
\begin{tabular}{lrrr}
\toprule
Seed & $\Delta Y$ (s/departed) & $\Delta R_{\mathrm{cmp}}$ & $\Delta U$ (passengers) \\
\midrule
6201 & $-87.25$ & $+0.00457$ & $-75.43$ \\
6202 & $-97.44$ & $+0.00485$ & $-80.67$ \\
6203 & $-94.28$ & $+0.00478$ & $-79.63$ \\
6204 & $-93.44$ & $+0.00465$ & $-76.97$ \\
6205 & $-88.21$ & $+0.00444$ & $-73.40$ \\
\bottomrule
\end{tabular}
\end{table}

\subsection{Formal evidence package}
\label{app:evidence}

The formal package records 45 physical behaviors, 30 fresh environment cells,
ten complete block manifests, and 1,350 result records. The registered analyzer
reads only endpoint results and implements the exact sign tests and gates in
\Cref{sec:exp:statistics}. A second analyzer independently recomputes result
counts, candidate--parent means, demand and block differences, win counts, and
gate outcomes. A separate closeout verifies source, protocol, environment,
method, scheduler, and file-manifest bindings.

The central machine-readable files are:
\begin{itemize}[leftmargin=*,nosep]
    \item \path{formal_confirmation_analysis_v1.json};
    \item \path{formal_confirmation_analysis_verification_v1.json};
    \item \path{evaluation_closeout_v1.json}; and
    \item \path{evaluation_closeout_verification_v1.json}.
\end{itemize}
\begin{samepage}
They are stored in the following directory:
\begin{quote}
\small
\texttt{reproducibility/}\\
\texttt{h2oplus\_v20\_formal\_confirmation\_results/}
\end{quote}
The synchronized package intentionally contains no checkpoint, optimizer,
dataset, replay, or remote-workspace copy.
\end{samepage}
